\documentclass[a0,portrait]{hogent-poster}

% Info over de opleiding
\course{Graduaatsproef}
\studyprogramme{Graduaat in het Programmeren}
\academicyear{2025-2026}
\institution{Hogeschool Gent, Valentin Vaerwyckweg 1, 9000 Gent}

% Info over de bachelorproef
\title{Automatisering van Release-verificatie binnen de Amaron Releasecyclus.}
\author{De Pus Jonas}
\email{jonas.depus@student.hogent.be}
\supervisor{Willem Delvaux (Hogeschool)}
\cosupervisor{Koen Casier (Bedrijf)}

\specialisation{Programmeren}
\keywords{DevOps, CI/CD, Release Management, API, Automatisering}
\projectrepo{https://github.com/JonasDeP/gradproef-2026.git}

\begin{document}

\raggedbottom 

\maketitle

\begin{abstract}
Binnen Amaron zorgt de handmatige controle van circa 40 softwarecomponenten voorafgaand aan elke release voor veel tijdverlies, fouten en frustraties. Het ontbreken van een centraal overzicht leidt tot vertragingen en miscommunicatie tussen developers en het DevOps-team.

Om dit op te lossen, is een geautomatiseerd verificatieplatform ontwikkeld. Developers registreren hun afgewerkte componenten nu zelf, waarna API-koppelingen de status automatisch en onafhankelijk verifiëren. Deze data wordt gebundeld in één centraal dashboard, waardoor het DevOps-team in één oogopslag ziet welke componenten klaar zijn voor de uitrol. Dit platform elimineert tijdrovende handmatige checks, reduceert de foutgevoeligheid drastisch en zorgt voor kortere doorlooptijden. Het resultaat is een vlottere samenwerking, minder frustraties en een schaalbare oplossing voor toekomstige groei.
\end{abstract}

\begin{multicols}{2} 

\section{Introductie}

Binnen de moderne softwareontwikkeling is een efficiënt en betrouwbaar releaseproces cruciaal om kwalitatieve oplossingen snel tot bij de eindgebruiker te krijgen. Bij het softwarebedrijf Amaron wordt gewerkt met een vaste releasecyclus van twee tot vier weken. Het bedrijf is opgebouwd uit een ecosysteem van inmiddels circa 40 verschillende componenten en modules, die gebundeld in de hoofdrelease worden uitgerold.

Naarmate dit ecosysteem is gegroeid, is de complexiteit van het releasebeheer evenredig toegenomen. Waar in het verleden het overzicht eenvoudig te behouden was, vereist de huidige schaal een meer gestructureerde aanpak. Deze graduaatsproef richt zich op de optimalisatie en automatisering van dit specifieke onderdeel binnen de releasecyclus van Amaron.

\section{Experimenten}

Om de betrouwbaarheid en gebruiksvriendelijkheid van het platform te garanderen, is het systeem in drie fasen getest:

\begin{itemize}
  \setlength{\itemsep}{0.5em}
  \setlength{\parskip}{0pt}
  \item \textbf{Technische validatie:} De ontwikkelde API-koppelingen zijn uitgebreid getest om een foutloze dataverwerking en communicatie met de achterliggende systemen te verzekeren.
  \item \textbf{End-to-End simulatie:} In een testomgeving is een volledige releasecyclus gesimuleerd om te verifiëren of geregistreerde componenten correct en realtime op het centrale dashboard verschijnen.
  \item \textbf{Gebruikerstesten:} Samen met de developers is het registratieformulier in de praktijk getest en geoptimaliseerd, zodat de flow intuïtief aanvoelt en een vlotte adoptie binnen het team gegarandeerd is.
\end{itemize}

\section{Conclusies}

Dit project toont aan dat het handmatige en foutgevoelige releaseproces bij Amaron succesvol geautomatiseerd kan worden. Door een workflow te ontwikkelen waarbij developers hun afgewerkte componenten registreren en API-koppelingen de verificatie overnemen, krijgt het DevOps-team via één centraal dashboard direct een realtime overzicht. Deze oplossing reduceert niet alleen het tijdsverlies en de menselijke foutgevoeligheid, maar neemt ook de aanzienlijke communicatie-overhead en bijbehorende frustraties op de werkvloer weg.

\section{Toekomstig onderzoek}

Hoewel het platform momenteel technisch succesvol is, blijft het afhankelijk van een menselijke handeling aan de start van de flow. Om dit in de toekomst verder te optimaliseren, worden de volgende stappen aanbevolen:

\begin{itemize}
  \setlength{\itemsep}{0.5em}
  \setlength{\parskip}{0pt}
  \item \textbf{Volledige CI/CD integratie:} De initiële registratiestap wegnemen door automatische triggers in te bouwen bij een specifieke Git-merge.
  \item \textbf{Extra API-controles:} Het automatisch ophalen van security scans en testrapporten toevoegen aan de verificatieflow.
  \item \textbf{Gefaseerde uitrol:} Het systeem stapsgewijs implementeren voor alle 40+ componenten, inclusief training voor een optimale adoptie binnen het team.
\end{itemize}

\section{Sectie met figuur}
\begin{center}
  \captionsetup{type=figure}
  \includegraphics[width=0.9\linewidth]{form}
  \captionof{figure}{Dit voorbeeld toont de interface van een formulier, voorzien van diverse gegevensvelden voor zowel datavisualisatie als gebruikersinvoer, aangevuld met actieknoppen voor navigatie en functie-uitvoering.}
\end{center}

\begin{center}
  \captionsetup{type=figure}
  \includegraphics[width=1.0\linewidth]{id}
  \captionof{figure}{Op de bijgevoegde screenshot zie je de visuele weergave van de interaction designer. Deze opzet illustreert mooi hoe de verschillende elementen met elkaar verbonden zijn om complexe handelingen op de achtergrond te stroomlijnen.}
\end{center}

\vspace*{\fill}

\end{multicols}
\end{document}